\documentclass{article}

% ============================================================================
% F.A.B.L.E. — NeurIPS 2026 style preprint.
% Drop-in: once you add neurips_2026.sty to this folder, comment the \usepackage
% fallbacks below and uncomment the neurips line.
%   \usepackage[preprint]{neurips_2026}
% For anonymous submission:  \usepackage{neurips_2026}
% For a workshop:            \usepackage[dblblindworkshop]{neurips_2026}
%                            \workshoptitle{Workshop Name Here}
% ============================================================================

\usepackage[utf8]{inputenc}
\usepackage[T1]{fontenc}
\usepackage{hyperref}
\usepackage{url}
\usepackage{booktabs}
\usepackage{amsfonts}
\usepackage{amsmath}
\usepackage{amssymb}
\usepackage{nicefrac}
\usepackage{microtype}
\usepackage{xcolor}
\usepackage{graphicx}
\usepackage{enumitem}
\usepackage{multirow}
\usepackage{array}
\usepackage{hyphenat}
\usepackage{tikz}
\usetikzlibrary{arrows.meta, positioning, fit, backgrounds, calc}

% Lightweight margins so the article fallback resembles the NeurIPS column width.
\usepackage[margin=1in]{geometry}

% ---- TikZ node styles for the architecture diagrams ------------------------
\tikzset{
  stage/.style   = {rectangle, rounded corners=2pt, draw=black!70, fill=blue!6,
                    minimum height=8mm, minimum width=20mm, align=center, font=\small},
  agent/.style   = {rectangle, rounded corners=2pt, draw=black!70, fill=purple!8,
                    minimum height=8mm, minimum width=20mm, align=center, font=\small},
  gov/.style     = {rectangle, rounded corners=2pt, draw=black!70, fill=green!8,
                    minimum height=8mm, minimum width=20mm, align=center, font=\small},
  store/.style   = {cylinder, shape border rotate=90, aspect=0.3, draw=black!70,
                    fill=orange!10, minimum height=11mm, minimum width=16mm,
                    align=center, font=\scriptsize},
  io/.style      = {rectangle, draw=black!70, fill=black!4, minimum height=8mm,
                    align=center, font=\small},
  flow/.style    = {-{Stealth[length=2mm]}, thick},
  dflow/.style   = {-{Stealth[length=2mm]}, thick, dashed, draw=black!55},
}

\title{F.A.B.L.E.: A Framework for Adversarial Benchmarking and \\ Logic Evaluation in Privacy-Aware Multi-Agent LLM Systems}

\author{%
  Pratik Balaji \\
  Temple University \\
  Philadelphia, PA \\
  \texttt{balajipratik8@gmail.com} \\
}
\date{}

\begin{document}

\maketitle

\begin{abstract}
We present F.A.B.L.E., a privacy-aware framework for studying how multiple language-model
agents cooperate, disagree, and converge under explicit governance constraints. Rather than
treating a model's answer as a single opaque string, F.A.B.L.E. runs two contrasting
orchestration protocols over the same evaluation harness: a lightweight three-stage
\emph{standard} pipeline (Analyst $\rightarrow$ Critic $\rightarrow$ Synthesizer) and a
six-stage \emph{adversarial} pipeline (Planner $\rightarrow$ Actor $\rightarrow$ Critic
$\rightarrow$ Validator $\rightarrow$ Refiner $\rightarrow$ Judge) with bounded
refine-and-rejudge rounds. Every request passes through a fail-closed governance layer
(regex$+$LLM PII redaction with an optional Microsoft Presidio sidecar, two-layer content
guardrails, and AES-256-GCM credential encryption), is grounded by a corrective
retrieval layer (CRAG-style relevance grading over a persisted FAISS corpus and per-user
semantic memory), and is scored by a five-dimension rubric that yields an auditable verdict.
We describe the architecture and its governance and retrieval design, and report two rounds of
head-to-head evaluation: an initial ten-run batch on largely-memorized prompts, and a
follow-up 60-case suite (20 standard, 20 adversarial, 10 Monte Carlo) redesigned so every
prompt has a genuine flaw to catch --- interacting constraints, a planted false premise, or a
claim that must be checked against a seeded retrieval corpus --- run under three orchestrator
backends (native asyncio, LangGraph, LangChain). On the discriminative suite the adversarial
protocol's mean score is 15--18 points below the standard committee's (72--74\% vs.\ 87--89\%,
consistent across all three orchestrators) at roughly $1.3\times$ the latency, the opposite
framing from the original batch and, we argue, the more informative one. Diagnosing this gap
surfaced a measurement artifact in our own harness: the Judge's system prompt forces an ACCEPT
verdict on the final refine round regardless of quality, so 100\% (60/60) of adversarial runs
in this dataset accepted by construction rather than genuine judgment, making the numeric score
the only real signal. We report this finding, its fix, and its implications alongside the
headline numbers. We position the system as a reproducible platform for measuring
\emph{disagreement and synthesis}, not only final-answer accuracy, and release the
orchestration code, governance layer, and evaluation harness as an integrated artifact.
\end{abstract}

\section{Introduction}
Language models are now deployed as assistants, evaluators, and components of autonomous
workflows, yet they still produce confidently-worded answers whose intermediate reasoning is
brittle, contradictory, or unauditable. Multi-agent arrangements---debate, critique, role
specialization---can add useful diversity, but the gains are sensitive to protocol design and
notoriously hard to inspect: majority agreement is not correctness when models share
correlated errors. At the same time, any system that touches user prompts inherits privacy
and governance obligations that are usually bolted on after the fact rather than designed in.

F.A.B.L.E.\ was built to study both problems with one apparatus. The premise is that an AI
answer should be decomposed into inspectable stages that can be benchmarked individually.
A request is sanitized, screened by safety guardrails, grounded by graded retrieval, routed
through an explicit multi-stage agent pipeline, scored against a rubric, and returned with a
verdict---after the original entities masked during sanitization are reinjected so the user
still reads a natural answer. Crucially, the framework ships two pipelines that share this
spine: a fast standard committee and a slower adversarial protocol in which a Critic and
Validator actively attack the Actor's draft and a Judge arbitrates after bounded refinement
rounds. Running both over the same harness turns ``does adversarial interaction help?'' into
a measurable question rather than a rhetorical one.

This paper reflects the system as actually implemented---the governance layer, both
pipelines, the retrieval stack, and the evaluation harness are running code, not a roadmap.
We are deliberate, however, about the scope of the empirical claims: the results here combine
an initial ten-run batch (five prompts per mode) with a follow-up 60-case, three-orchestrator
suite built to be discriminative rather than memorizable. Both are intended to demonstrate the
harness, surface the trade-off, and---in the second batch's case---surface a genuine artifact
in the evaluation methodology itself, not to establish a large-scale benchmark with statistical
guarantees.

\paragraph{Contributions.}
\begin{enumerate}[leftmargin=1.6em, itemsep=2pt]
  \item A modular, privacy-aware architecture for multi-agent LLM evaluation in which
        sanitization, guardrails, retrieval, orchestration, and scoring are separable,
        inspectable stages (Section~\ref{sec:arch}).
  \item Two contrasting orchestration protocols---a three-stage standard committee and a
        six-stage adversarial pipeline with bounded refine-and-rejudge rounds---sharing one
        governance and evaluation spine (Sections~\ref{sec:standard}--\ref{sec:adversarial}).
  \item A governance layer treated as a first-class research variable: fail-closed PII
        redaction with reinjection, two-layer guardrails, and AAD-bound credential
        encryption, validated by a 40-finding internal security audit (Section~\ref{sec:gov}).
  \item A corrective-retrieval stack (CRAG-style grading over FAISS $+$ semantic memory) and
        a golden-case reasoning cache for token/latency reduction (Section~\ref{sec:retrieval}).
  \item A reproducible evaluation harness; a preliminary ten-run, two-mode comparison and a
        follow-up 60-case, three-orchestrator, discriminative-prompt suite (verdict, rubric
        score, latency); and a diagnosed measurement artifact---the Judge's forced final-round
        ACCEPT makes the adversarial accept rate 100\% by construction under a two-round
        cap---reported alongside the headline numbers rather than smoothed over
        (Section~\ref{sec:eval}).
\end{enumerate}

\section{Related Work}
\textbf{Multi-agent reasoning.} A growing line of work has language-model agents debate,
critique, or specialize to improve answer quality, often reporting gains over single-model
baselines. The recurring caveat is that improvements depend heavily on the interaction
protocol and that consensus can mask correlated error. F.A.B.L.E.\ contributes a controlled
substrate in which a critique-driven adversarial protocol and a static committee are run side
by side under identical governance and scoring, and in which disagreement is recorded rather
than averaged away.

\textbf{LLM evaluation.} Many evaluations emphasize task accuracy or pairwise preference and
underweight intermediate structure, contradiction, or deployment-realistic governance.
F.A.B.L.E.\ complements these by scoring along five rubric dimensions and emitting a verdict
with rationale, and by making the orchestration cost (latency, round count) a reported metric
rather than an implementation detail.

\textbf{Retrieval-augmented generation.} Vector retrieval improves grounding but raises
questions about evidence trust and memory retention. F.A.B.L.E.\ adopts a corrective-RAG
posture: retrieved passages are graded for relevance, the query is rewritten and re-retrieved
when results are weak, and retrieved text is explicitly framed to agents as untrusted data
rather than instructions, mitigating indirect prompt injection.

\textbf{Privacy-aware NLP.} Input sanitization and PII redaction are usually production
concerns. Here they are part of the method: redaction changes the evidence visible to later
stages, so it is modeled and reported, with a reinjection step that restores user-facing
fidelity without persisting raw identifiers.

\section{System Architecture}
\label{sec:arch}
F.A.B.L.E.\ is a five-layer system: a Next.js/React web client; a Python (FastAPI)
orchestration core; a provider-agnostic model-access layer (OpenRouter-compatible, also
usable against any OpenAI-compatible endpoint and a no-API-key local embedder); a memory and
storage layer (FAISS for the document corpus, Supabase/pgvector for per-user semantic memory,
Amazon S3 for artifacts); and an evaluation layer (rubric scoring, verdict derivation, and a
live latency harness). The orchestration core exposes both pipelines behind a shared request
path so that sanitization, guardrails, retrieval, and reinjection are identical across modes
and only the agent topology differs.

\subsection{Standard pipeline}
\label{sec:standard}
The standard pipeline (Figure~\ref{fig:standard}) is a three-stage committee. After
governance and retrieval, an \emph{Analyst} produces a structured analysis, a \emph{Critic}
identifies gaps and unsupported claims, and a \emph{Synthesizer} produces the final answer.
The run is scored by the rubric and a PASS/WARN/FAIL verdict is derived from the rubric mean.

\begin{figure}[t]
\centering
\begin{tikzpicture}[node distance=6mm and 9mm]
  \node[io]    (in)   {User\\prompt};
  \node[gov, right=of in]    (pii)  {PII\\redact};
  \node[gov, right=of pii]   (guard){Guardrail\\pre-check};
  \node[stage, right=of guard](rag) {Agentic\\RAG};
  \node[agent, below=10mm of rag] (an) {Analyst};
  \node[agent, left=of an]  (cr) {Critic};
  \node[agent, left=of cr]  (sy) {Synthesizer};
  \node[stage, left=of sy]  (rub){Rubric\\$+$ verdict};
  \node[gov, below=10mm of rub] (re){PII\\reinject};
  \node[io, right=of re] (out){Final\\answer};

  \draw[flow] (in) -- (pii);
  \draw[flow] (pii) -- (guard);
  \draw[flow] (guard) -- (rag);
  \draw[flow] (rag) -- (an);
  \draw[flow] (an) -- (cr);
  \draw[flow] (cr) -- (sy);
  \draw[flow] (sy) -- (rub);
  \draw[flow] (rub) -- (re);
  \draw[flow] (re) -- (out);

  % side stores
  \node[store, above=7mm of rag] (corpus){FAISS\\corpus};
  \node[store, left=10mm of corpus] (mem){pgvector\\memory};
  \draw[dflow] (corpus) -- (rag);
  \draw[dflow] (mem) -- (rag);
  \node[store, below=7mm of sy] (kg){Knowledge\\graph};
  \draw[dflow] (rub) |- (kg);
\end{tikzpicture}
\caption{Standard pipeline. Green = governance, blue = system stages, purple = agents,
orange = stores. Retrieval grounds the Analyst; the completed run feeds the knowledge graph.}
\label{fig:standard}
\end{figure}

\subsection{Adversarial pipeline}
\label{sec:adversarial}
The adversarial pipeline (Figure~\ref{fig:adversarial}) shares the governance and retrieval
front-end but replaces the committee with a six-role protocol. A \emph{Planner} decomposes
the task; an \emph{Actor} drafts an answer; a \emph{Critic} attacks it for flaws; a
\emph{Validator} checks factual grounding and consistency against the (untrusted) retrieved
context; a \emph{Refiner} emits a structured fix specification; and the Actor revises. This
refine loop repeats up to a hard-capped number of rounds (default two), after which a
\emph{Judge} renders an ACCEPT/REJECT verdict with score and rationale. Role-to-model
assignment is deliberately heterogeneous to avoid correlated error across the Critic and
Refiner.

\begin{figure}[t]
\centering
\begin{tikzpicture}[node distance=5mm and 7mm]
  \node[io]    (in)   {User\\prompt};
  \node[gov, right=of in]   (pii)  {PII\\redact};
  \node[gov, right=of pii]  (guard){Guardrail};
  \node[stage, right=of guard](rag){Agentic\\RAG};

  \node[agent, right=of rag] (pl){Planner};
  \node[agent, below=9mm of pl] (ac){Actor};
  \node[agent, left=of ac] (cri){Critic};
  \node[agent, left=of cri](va){Validator};
  \node[agent, left=of va] (rf){Refiner};
  \node[agent, below=9mm of rf] (ju){Judge};
  \node[gov, right=of ju] (re){PII\\reinject};
  \node[io, right=of re] (out){Final\\answer\\$+$ verdict};

  \draw[flow] (in) -- (pii);
  \draw[flow] (pii) -- (guard);
  \draw[flow] (guard) -- (rag);
  \draw[flow] (rag) -- (pl);
  \draw[flow] (pl) -- (ac);
  \draw[flow] (ac) -- (cri);
  \draw[flow] (cri) -- (va);
  \draw[flow] (va) -- (rf);
  % refine loop back to Actor
  \draw[flow, dashed, draw=red!60] (rf) to[out=70,in=110] node[above,font=\scriptsize]{refine $\leq 2$ rounds} (ac);
  \draw[flow] (rf) -- (ju);
  \draw[flow] (ju) -- (re);
  \draw[flow] (re) -- (out);

  \node[store, above=6mm of rag] (corpus){FAISS\\corpus};
  \node[store, left=8mm of corpus] (mem){pgvector\\memory};
  \draw[dflow] (corpus) -- (rag);
  \draw[dflow] (mem) -- (rag);
\end{tikzpicture}
\caption{Adversarial pipeline. The Critic and Validator attack the Actor's draft; the Refiner
specifies fixes and the Actor revises for up to two rounds (red loop); the Judge arbitrates.
The golden-case cache is consulted and populated by the standard pipeline only
(Section~\ref{sec:standard}); the adversarial path always executes in full.}
\label{fig:adversarial}
\end{figure}

\section{Governance and Privacy}
\label{sec:gov}
Governance is part of the method because it alters the evidence later stages see.

\textbf{PII redaction with reinjection.} Each prompt is redacted before any model call: a
regex layer catches structured identifiers (email, phone, SSN, credit-card with Luhn check,
IBAN, IP, API keys) and an optional layer---either a small LLM extractor or a Microsoft
Presidio sidecar reachable over HTTP---catches names, locations, and organizations. Detected
spans are replaced by nonce-keyed placeholders; the real values are restored by a reinjection
pass on the outgoing answer so the user sees a natural response while raw identifiers never
enter logs or memory. The pipeline is \emph{fail-closed}: if redaction cannot complete,
orchestration halts.

\textbf{Two-layer content guardrails.} A microsecond rule layer (prompt-injection patterns,
credential-exfiltration patterns, a hate-speech block list, and Unicode-normalized matching
to defeat homoglyph evasion) runs before any model call; a cheap LLM classifier handles
nuanced cases and \emph{fails to ``warn,'' never to ``allow.''} Long inputs are chunked so an
injection hidden past a truncation boundary is still screened.

\textbf{Secrets.} User-supplied provider credentials are encrypted with AES-256-GCM bound to
the owning identity via additional authenticated data, so ciphertext is not portable across
users; a versioned key scheme supports rotation.

\textbf{Audit.} An internal security review enumerated 40 findings across authentication,
tenant isolation, CSRF/CORS, PII flow, secrets, guardrails, RAG/SSRF, and deployment; all
non-accepted findings were remediated and covered by regression tests. We treat this audit as
part of the artifact's reproducibility rather than as a compliance footnote.

\section{Retrieval and Memory}
\label{sec:retrieval}
\textbf{Corrective retrieval (CRAG-lite).} Retrieval pools candidates from a persisted FAISS
document corpus and, in multi-user mode, per-user pgvector memory. An LLM grades each
candidate for relevance; if too few survive, the query is rewritten and retrieval repeats up
to a bounded number of hops. The surviving passages are injected under an explicit
``untrusted data, not instructions'' framing. This makes the previously-unused document
corpus an active part of inference and reduces context pollution relative to single-shot
top-$k$.

\textbf{Golden-case reasoning cache.} A completed run that clears a quality bar (rubric mean
above threshold and a passing verdict) is promoted to a ``golden case'' stored with its
trajectory. A new prompt that is highly similar to a golden case can be answered by adapting
the stored answer (with a cheap re-check gate) or by warm-starting the pipeline with the
stored trajectory, trading a small re-check cost for a large reduction in tokens and latency.

\textbf{Adaptive routing.} A lightweight learned router (an extreme-learning-machine
meta-scorer over run embeddings) and a knowledge graph of past runs provide model-selection
hints and prior context; these are advisory and fall back to heuristics when untrained.

\section{Evaluation}
\label{sec:eval}
\textbf{Setup.} We use a live harness that issues five fixed prompts spanning code,
arithmetic reasoning, finance, factual summarization, and creative naming, each through both
the standard (\texttt{/run}) and adversarial (\texttt{/adversarial-run}) endpoints---ten runs
total. Each run records the verdict, the five-dimension rubric mean, the completed round count
(adversarial), and end-to-end latency. Quality is scored by the rubric (accuracy, depth,
clarity, actionability, coverage); the verdict is PASS/WARN/FAIL for the standard mode
(thresholded on the rubric mean) and ACCEPT/REJECT for the adversarial mode (the Judge's
decision). Configuration: run-level summaries only, Critic $=$ \texttt{claude-3.5-haiku},
Refiner $=$ \texttt{gpt-4o-mini}, \texttt{max\_rounds} $=2$, 45\,s per-call timeout. Batch run
2026-06-13.

\textbf{Standard mode.} Table~\ref{tab:standard} reports the five standard runs: mean rubric
score 82\%, mean latency 32.8\,s, pass rate 4/5. The single non-pass was the open-ended
creative prompt (WARN, 60\%), where the rubric penalized coverage and actionability.

\textbf{Token cost.} Per-call token usage is captured in \texttt{ModelResponse.usage}
(prompt + completion token counts) and priced via a model price table in
\texttt{backend/core/cost.py}. Estimated cost per run at current OpenRouter rates:
standard $\approx\$0.003$, adversarial $\approx\$0.030$, Monte Carlo (4 variants)
$\approx\$0.045$. The 50 pending cases of the 60-case suite there\-fore cost
$\approx\$1.10$ ($20\times0.003 + 20\times0.030 + 10\times0.045$), and we budget
\$1--\$2.50 to absorb multi-round variance. Aggregated costs are surfaced in the live
benchmark dashboard and reported per run in the runner output; we note that the runner's
own total is a lower bound, since it prices a run's aggregate token usage against the
single model identifier the endpoint returns while an adversarial run spans several
models.

\begin{table}[t]
\centering
\caption{Standard pipeline (Analyst $\rightarrow$ Critic $\rightarrow$ Synthesizer), 5 runs.}
\label{tab:standard}
\begin{tabular}{@{}clccr@{}}
\toprule
Run & Prompt & Verdict & Rubric score & Latency (s) \\
\midrule
S1 & code      & PASS & 87\% & 39.9 \\
S2 & reasoning & PASS & 90\% & 22.0 \\
S3 & finance   & PASS & 92\% & 45.2 \\
S4 & factual   & PASS & 82\% & 29.0 \\
S5 & creative  & WARN & 60\% & 27.7 \\
\midrule
\multicolumn{3}{r}{\textbf{Mean}} & \textbf{82\%} & \textbf{32.8} \\
\bottomrule
\end{tabular}
\end{table}

\textbf{Adversarial mode.} Table~\ref{tab:adversarial} reports the five adversarial runs:
mean rubric score 80\%, mean latency 72.3\,s, accept rate 5/5, every run using both refine
rounds. Notably the creative prompt that only WARNed in the standard committee was driven to
an ACCEPT (82\%) by the critique-and-refine loop, and the finance run---whose Judge output was
truncated by the token budget---was correctly recovered as ACCEPT/75\% by a field-salvage
parser rather than being silently rejected (see Section~\ref{sec:findings}).

\begin{table}[t]
\centering
\caption{Adversarial pipeline (Planner $\rightarrow$ Actor $\rightarrow$ Critic $\rightarrow$
Validator $\rightarrow$ Refiner $\rightarrow$ Judge), 5 runs.}
\label{tab:adversarial}
\begin{tabular}{@{}clcccr@{}}
\toprule
Run & Prompt & Verdict & Judge score & Rounds & Latency (s) \\
\midrule
A1 & code      & ACCEPT & 85\% & 2/2 & 78.9 \\
A2 & reasoning & ACCEPT & 85\% & 2/2 & 61.6 \\
A3 & finance   & ACCEPT & 75\% & 2/2 & 90.6 \\
A4 & factual   & ACCEPT & 75\% & 2/2 & 68.1 \\
A5 & creative  & ACCEPT & 82\% & 2/2 & 62.4 \\
\midrule
\multicolumn{4}{r}{\textbf{Mean}} & & \textbf{72.3} \\
\bottomrule
\end{tabular}
\end{table}

\textbf{Head-to-head.} Table~\ref{tab:headtohead} summarizes the trade-off. The adversarial
protocol raised the worst-case prompt from WARN to ACCEPT and reached a 5/5 accept rate, but
at $\sim$2.2$\times$ the latency of the standard committee. Mean rubric scores are close
(82\% vs.\ 80\%), which is itself informative: the adversarial protocol's value in this batch
showed up as \emph{reduced variance and eliminated failures} rather than a higher mean,
consistent with the framing that critique-driven interaction trades cost for reliability.
We stress that $n=5$ prompts per mode in a single batch is preliminary; we report it to
demonstrate the harness and the measured trade-off, not as a statistically powered result.
\textbf{As Section~\ref{sec:findings} details, this 5/5 accept rate is not evidence of quality
on its own} --- every run here used both refine rounds, and the Judge's prompt forces an ACCEPT
verdict on the final round regardless of score.

\begin{table}[t]
\centering
\caption{Head-to-head summary (5 prompts per mode, single batch).}
\label{tab:headtohead}
\begin{tabular}{@{}lccc@{}}
\toprule
Mode & Pass/Accept rate & Mean score & Mean latency (s) \\
\midrule
Standard    & 4/5 & 82\% & 32.8 \\
Adversarial & 5/5 & 80\% & 72.3 \\
\bottomrule
\end{tabular}
\end{table}

\subsection{60-case discriminative re-evaluation}
\label{sec:sixty-case}
The five-prompt batch above uses canonical, largely memorized items (e.g.\ the bat-and-ball
problem) with no planted flaw for the Critic or Validator to catch, which likely explains why
the adversarial mean did not exceed the standard mean. We re-ran the harness on a redesigned
20-prompt set (4 per category $\times$ 5 categories: code, reasoning, factual, document QA,
writing) built so every prompt has something to critique: interacting constraints an obvious
solution violates (e.g.\ ``$O(n)$ time \emph{and} $O(1)$ extra space \emph{and} handle
duplicates''), a planted false premise, or a claim that must be checked against a seeded
retrieval corpus rather than parametric memory. We also seeded a dedicated corpus (EU AI Act,
retrieval-augmented generation, nearest-neighbor search, Roth IRA) so the document-QA category
actually exercises the CRAG-lite retrieval path instead of degenerating into recall. Each
configuration ran all 60 cases (20 standard, 20 adversarial, 10 Monte Carlo) under three
orchestrator backends (native \texttt{asyncio}, LangGraph, LangChain) to test whether they
execute the same pipeline within noise.

\begin{table}[t]
\centering
\caption{60-case suite, redesigned discriminative prompts, one orchestrator per row.}
\label{tab:sixtycase}
\begin{tabular}{@{}lcc@{}}
\toprule
Orchestrator & Std.\ score/latency/pass & Adv.\ score/latency/accept \\
\midrule
asyncio   & 89\% / 86.4s / 20\slash20 & 72\% / 105.1s / 20\slash20 \\
langgraph & 89\% / 82.5s / 20\slash20 & 74\% / 114.7s / 20\slash20 \\
langchain & 87\% / 88.2s / 19\slash20 & 71\% / 115.8s / 20\slash20 \\
\bottomrule
\end{tabular}
\end{table}

Unlike the five-prompt batch, the adversarial mean here is clearly below the standard mean
(72--74\% vs.\ 87--89\%) with a wide score spread (std.\ dev.\ 12\%, range 35--92\% across all
60 adversarial runs pooled), confirming the redesigned prompts give the Critic/Validator/Judge
something to actually discriminate on. The three orchestrators agree within noise on both
modes --- the one exception is \texttt{langchain}'s standard pass, which WARNed on one docqa
case that PASSed at 87--90\% under the other two backends, discussed alongside that case's
adversarial instability in Section~\ref{sec:findings}. Per-category adversarial means (pooled
across all three orchestrator runs, $n=12$/category) rank: reasoning 85\%, factual 75\%,
writing 73\%, code 71\%, and document QA 58\% (range 35--72\%) --- the corpus-grounding
constraint is the hardest for the Refiner to fully satisfy within the two-round budget.

\subsection{Root-cause findings during evaluation}
\label{sec:findings}
The harness doubled as a debugging instrument. Three findings are worth recording because they
affect any reproduction. First, an early adversarial batch hit the client timeout: two roles
were configured with a model ID that 404s on the provider, forcing a failed call plus fallback
retry every round, compounded by inline per-agent summarization; fixing the IDs and moving to
run-level summaries cut roughly eleven calls per run and brought peak latency to 90.6\,s.
Second, two runs initially reported REJECT/0\% when the Judge had in fact returned ACCEPT---the
JSON verdict was truncated by the token budget and the parser fell through to its reject
default. Raising the Judge budget and adding a field-level salvage parser recovered the valid
acceptances and moved the adversarial accept rate from 3/5 to 5/5.

Third, and most consequential: investigating why the 60-case adversarial mean sat well below
the standard mean (Section~\ref{sec:sixty-case}) surfaced a measurement artifact rather than
only a capability gap. The Judge's system prompt instructs it to force an ACCEPT verdict on the
final round \emph{regardless of quality}, so the pipeline always terminates rather than looping
indefinitely. Checking \texttt{rounds\_completed} across all 60 adversarial-mode records from
the three orchestrator runs (and, retroactively, the five-run legacy batch above) shows
\textbf{100\% (60/60) hit the round cap with verdict ACCEPT} --- not one run in the dataset
organically accepted before the final round. \emph{The ACCEPT/REJECT verdict is therefore not a
discriminating signal under the current two-round cap; it is 100\% by construction, and the
numeric Judge score is the only quality signal that varies.} This is corroborated qualitatively:
pulling the (truncated) rationale text for the three lowest-scoring/most volatile cases across
all three orchestrator runs shows the Judge correctly identifying the same planted violation at
round 2 that it flagged at round 1 in every instance --- a complexity-constraint violation for a
code prompt (scores 45--65\%) and unsourced or fabricated claims against the retrieval corpus
for two document-QA prompts (scores 35--72\%). The harness is working as designed; the Refiner
simply does not always resolve a flagged issue within one revision. We added a
\texttt{forced\_accept} field to the API response (true when
\texttt{judge\_verdict==ACCEPT} and \texttt{rounds\_completed==max\_rounds}) so future runs
report this directly rather than requiring retroactive log analysis. We did not have budget
remaining to test whether raising \texttt{adversarial\_max\_rounds} above 2 reduces the
forced-accept rate, nor were full per-round transcripts persisted for this batch to confirm
\emph{why} the Refiner's revision failed rather than only that the Judge still flagged the same
issue --- both are noted as immediate future work in Section~\ref{sec:limitations}.

\section{Implementation Status}
The governance layer (PII redaction with reinjection, Presidio sidecar option, two-layer
guardrails, AAD-bound encryption), both orchestration pipelines, the corrective-retrieval and
golden-case layers, the rubric/verdict evaluator, the live latency harness, and the web client
are implemented and exercised by an automated test suite. Infrastructure for containerized and
Kubernetes deployment, plus an AWS (ECS/EKS) path, is defined as infrastructure-as-code.
Preliminary: the ten-run batch is a single seed; multi-user pgvector memory and the
managed-cloud deployments are validated by code review and mocked tests rather than a live
production run. We are explicit about this boundary so the empirical claims are not overread.
The adversarial pipeline additionally supports an optional run-level self-consistency
ensemble (\texttt{ADVERSARIAL\_ENSEMBLE\_SIZE}): $N$ debates are sampled independently and
reduced by majority vote over the final answer, falling back to highest Judge score when no
majority exists. Budget constraints limited live verification to $N=2$ on 3 cases rather than
the intended $N=5$ over the full suite: the reducer's fallback path (score tie-break when no
majority answer forms) fired correctly in all 3 cases, and picked the higher-scoring candidate
in the 2 non-tied cases, but $N=2$ rarely produces an exact-text majority, so the majority-vote
branch itself remains unverified on live data.

\section{Limitations}
\label{sec:limitations}
\textbf{Small evaluation.} Five prompts per mode in one batch is a demonstration, not a
benchmark; there are no confidence intervals, no human evaluation, and no comparison against
external multi-agent baselines on a standard dataset.
\textbf{Rubric and judge are model-based.} Both the rubric scorer and the adversarial Judge
are themselves LLMs and inherit their biases; verdicts should be read as structured estimates,
not ground truth.
\textbf{Provider dependence.} Although provider-agnostic in principle, runs depend on external
model availability and on the chosen role-to-model mapping.
\textbf{Privacy trade-off.} Redaction reduces exposure but can remove context useful for
reasoning; quantifying that effect is future work.
\textbf{Single-process caches.} The golden-case cache and per-identity concurrency limit are
in-process and do not yet coordinate across distributed replicas.
\textbf{Ensemble majority-vote branch unverified at scale.} The 3-case, $N=2$ probe
(Section~\ref{sec:findings}) exercised the reducer's fallback path only; the full $N=5$,
60-case ensemble benchmark remains future work pending budget.
\textbf{Adversarial ACCEPT rate is not a discriminating metric under the current round cap.}
Every adversarial run in the 60-case suite (and the original five-run batch) hit
\texttt{rounds\_completed==max\_rounds} with a forced ACCEPT verdict
(Section~\ref{sec:findings}); any reported ACCEPT/REJECT rate under
\texttt{adversarial\_max\_rounds}$=2$ should be read as 100\% by construction, not as evidence
of quality --- only the numeric Judge score discriminates. Whether a higher round cap lets the
Judge organically accept or reject, and whether the Refiner's revision genuinely fails on the
flagged cases or simply is not given enough attempts, are both open and require a rerun with
more rounds and persisted per-round transcripts.

\section{Conclusion}
F.A.B.L.E.\ is an attempt to make multi-agent LLM workflows inspectable, governable, and
comparable from the ground up. By sharing one governance, retrieval, and scoring spine across
a fast standard committee and a slower adversarial protocol, it turns ``does adversarial
interaction help?'' into a measured question. Our original ten-run batch used largely
memorized prompts and showed only a small accuracy-for-cost trade-off; redesigning the suite
to plant genuine flaws (Section~\ref{sec:sixty-case}) surfaced a real 15--18-point gap between
standard and adversarial scores, and, in the process, an artifact in our own harness --- the
Judge's forced-accept-on-final-round rule makes ACCEPT/REJECT uninformative under a two-round
cap, a finding we consider as valuable as the headline numbers. Three orchestrator backends
(native asyncio, LangGraph, LangChain) agree within noise on both pipelines. The contribution
is the platform and its framing---disagreement and synthesis as first-class, governance as a
research variable, and an evaluation harness that caught its own measurement bugs---together
with an honest account of what is built and what remains to be validated at scale.

\bibliographystyle{plainnat}
% \bibliography{references}

\appendix

\section{Reproducibility}
\textbf{Harness.} \texttt{scripts/eval\_runs.py} issues the five fixed prompts to
\texttt{/run} and \texttt{/adversarial-run} against a running backend, records latency and
verdict per run, and appends the results table to the project research log. \textbf{Config.}
Run-level summaries only; Critic \texttt{claude-3.5-haiku}; Refiner \texttt{gpt-4o-mini};
\texttt{adversarial\_max\_rounds}=2; 45\,s per-call timeout; client timeout 180\,s.
\textbf{Embeddings.} 384-dimensional; pluggable across a no-API-key local model
(\texttt{bge-small-en-v1.5}), OpenAI, or Gemini's OpenAI-compatible endpoint, keeping the
\texttt{vector(384)} schema and FAISS index aligned. \textbf{Determinism.} Model APIs are not
deterministic; absolute numbers will vary between batches, but the harness, prompts, and
scoring are fixed so the \emph{comparison} is reproducible.
\textbf{Declared initial state.} The system carries three adaptive stores that would
otherwise let one run influence the next: the golden-case cache, the knowledge-engine run
history, and the ELM meta-scorer's weights. Every reported run therefore starts from an
empty \texttt{data/knowledge/} (reset via \texttt{scripts/reset\_benchmark\_state.py}) and
executes with \texttt{GOLDEN\_CACHE\_ENABLED=false}. This matters specifically because the
standard pipeline is measured twice under different orchestrators over an identical prompt
set: with the cache live, the second pass would match its own cached entries at similarity
$\approx 1.0$ and return recycled answers rather than executing the pipeline.

\section{Rubric and Verdict}
Quality is scored on five dimensions---accuracy, depth, clarity, actionability, coverage. The
standard verdict thresholds the rubric mean into PASS/WARN/FAIL; the adversarial verdict is
the Judge's ACCEPT/REJECT with a numeric score and a one-paragraph rationale, parsed with a
salvage path that recovers truncated-but-valid verdicts.

\section{Security Audit Summary}
The internal review recorded 40 findings (6 critical, 14 high, 14 medium, 6 low). All
non-accepted findings were remediated---centralized tenant scoping, CSRF/CORS hardening,
fail-closed guardrails with Unicode normalization, AAD-bound credential encryption with key
rotation, SSRF protection on URL ingestion, per-identity concurrency limits, identity
revocation, and notebook-export sanitization---and covered by a dedicated regression test
module. Five low-impact items are documented as accepted risks with rationale.

\section{60 Preliminary Eval Test Cases}
\label{app:bench60}

Table~\ref{tab:bench60_finished} lists the ten completed Phase-14 runs. The remaining 50
cases (20 standard, 20 adversarial, 10 Monte Carlo) are defined in
\texttt{benchmarks/benchmark\_v1.yaml} and pending execution; placeholder rows indicate
\textsc{pending} verdicts. The full 60-case suite is the intended unit of reporting for
the camera-ready version.

\textbf{Structure.} Twenty shared prompts span five categories (four per category):
code generation, mathematical / logical reasoning, factual explanation, RAG/CRAG
document QA, and creative writing. Every shared prompt is issued to both the standard and
adversarial pipelines using the identical wording to enable a direct paired comparison.
The ten Monte Carlo prompts are the two most word-sensitive prompts per category,
selected for their expected high paraphrase-divergence signal (i.e.\ small wording
changes are expected to shift the response substantially).

\begin{table}[h!]
\centering
\small
\caption{Phase-14 finished runs (10/60). Full 60-case results pending; see
\texttt{benchmarks/BENCHMARK\_RESULTS.md}.}
\label{tab:bench60_finished}
\begin{tabular}{@{}lllccr@{}}
\toprule
Run & Category & Mode & Verdict & Score & Latency (s) \\
\midrule
S1 & code      & standard    & PASS   & 87\% & 39.9 \\
S2 & reasoning & standard    & PASS   & 90\% & 22.0 \\
S3 & factual   & standard    & PASS   & 92\% & 45.2 \\
S4 & factual   & standard    & PASS   & 82\% & 29.0 \\
S5 & writing   & standard    & WARN   & 60\% & 27.7 \\
A1 & code      & adversarial & ACCEPT & 85\% & 78.9 \\
A2 & reasoning & adversarial & ACCEPT & 85\% & 61.6 \\
A3 & factual   & adversarial & ACCEPT & 75\% & 90.6 \\
A4 & factual   & adversarial & ACCEPT & 75\% & 68.1 \\
A5 & writing   & adversarial & ACCEPT & 82\% & 62.4 \\
\midrule
\multicolumn{4}{l}{Standard mean / Adversarial mean} & 82\% / 80\% & 32.8 / 72.3 \\
\bottomrule
\end{tabular}
\end{table}

\textbf{Monte Carlo cases} (10): identifiers MC-C2, MC-C3 (code), MC-R2, MC-R3
(reasoning), MC-F3, MC-F4 (factual), MC-D1, MC-D3 (document QA), MC-W2, MC-W3
(writing). Each is issued with $n=4$ paraphrase variants; the Monte Carlo engine computes
pairwise cosine similarity across model responses to derive a consensus score and
divergence pairs (see Section~3).

\textbf{Reproducibility.} Run the pending 50 cases with:
\begin{verbatim}
python scripts/benchmark_v1.py
\end{verbatim}
Cost estimate: \$0.10--\$0.30 total. Results are written to
\texttt{data/benchmarks/results/} (raw JSON) and
\texttt{benchmarks/BENCHMARK\_RESULTS.md} (markdown). The benchmark dashboard at
\texttt{/dashboard} provides live mode analytics, token cost tracking, and an
OpenTelemetry trace waterfall per run. The Kaggle artifact (CSV + JSONL + reproducer
notebook) can be pushed from the dashboard's ``Export to Kaggle'' button.

\end{document}
